\scanpage{377}
\mspage{293}
% archive sheet number: 376

$\underline S(\pi)$. A fortiori $\underline S(\pi)$ est une sous-catégorie
str. pleine et \unsure{Karoubienne|karoubienne} de $M_\pi$.]

\deleted{\unsure{On aura|En outre}, puisque chaque} D'autre part, chaque $\underline S_\alpha$ étant stable \inserted{dans $M$} par
$\varinjlim$ gdes devant $\pi'_0$, il en est de $=$ de
leur intersection $\underline S$. A fortiori, celle-ci
est stable par $\varinjlim$ gdes dev.\textsuperscript{t} $\pi$, donc
le foncteur pl. fidèle canonique
$\operatorname{Ind}_\pi(\underline S(\pi))\unsure{\text{---}|\longrightarrow}M$ se factorise par $S$
en
\[
\operatorname{Ind}_\pi\underline S(\pi)\hookrightarrow\underline S,
\]
il faut prouver que ce foncteur est
\unsure{essent\textsuperscript{ellt}|essent\textsuperscript{llmt}} surjectif. Soit donc $x\in\deleted{\unsure{\operatorname{Ob}}}S=\operatorname{Ob}\underline S$,
tout revient à prouver que
\[
\varinjlim_{i\in\underline S(\pi)_{/x}}x_i\xrightarrow{\sim}x,\tag*{$(*)$}
\]
et que $\underline S(\pi)_{/x}$ est g\textsuperscript{r}de devant $\pi$.

\deleted{Il revient au même de} Notons que, comme $\pi'_0\geq\pi_0$, et \deleted{\unsure{on}}
$\pi^{\pi'_0}=\pi$, on aura $\pi^{\pi_0}=\pi$, donc $\pi$ est
cardinal utile pour $M$, donc \deleted{\unsure{on a}}
$I=M_{\pi/x}$ est g\textsuperscript{r}de dev.\textsuperscript{t} $x$, et on a
\[
\varinjlim_{i\in I=M_{\pi/x}}x_i\xrightarrow{\sim}x.\tag*{$(**)$}
\]
Donc, si l'on admet \inserted{que} $\underline S(\pi)_{/x}$ est \deleted{filtrante}\inserted{g\textsuperscript{r}de dev.\textsuperscript{t} $\pi$},
